%!TEX root = ../../main.tex
\subsection{Simulation Studies}

% OLD CONTENT
Corollary \ref{cor: three candidates} provides theoretical guidance on the choice of $\lamz$ to minimize the upper bound on excess risk. However, it does not guarantee that the same value of $\lamz$ minimizes the excess risk in practice. To investigate this, we conduct simulations to identify the values of $\lamz$ that minimize the risk of our algorithm on the test set (which is same as minimizing the excess risk).

Let 
\begin{equation*}
    \lamz^* := \arg\min_{\lamz} R(\whl),
\end{equation*}
 where $R(\cdot)$ denotes the risk evaluated on the test set, serving as a proxy for the population risk.

Our simulation study evaluates the behavior of $\lamz^*$ value, across two distinct scenarios: a collection of toy datasets and a collection of real-world datasets.\coderef{subsec:code_side_exp} 
\paragraph{Toy DataSet:}
We generate $100$ different toy distributions over $(X, Y, Z)$, each constructed as follows:

\begin{enumerate}
    \item The features $X$ are sampled from a Gaussian mixture model, whose parameters and mixing probabilities are randomly chosen.
    
    \item Then the weight vectors $u$ and $w$ are sampled from standard normal distribution.    
    Given fixed weights $u$, the binary variables $Z$ is sampled via logistic models:
    \[ Z \sim \sigma((1,X)^\top u).\]
    Once we obtain $Z$, for a fixed weight $w$, the binary variable $Y$ is sampled via a logistic model:
    \[ Y \sim \sigma((1,X,Z)^\top w).\]
    
\end{enumerate}

This process is repeated $100$ times, each with a new random choice of parameters, resulting in a diverse family of distributions. For each distribution, we generate $m=1000$ samples for the $(X, Z)$ dataset (used for learning $X \to Z$), $n=4000$ samples for the $(X, Y)$ dataset (used for learning $X \times Z \to Y$), and $t=20000$ samples for the test set $(X, Y, Z)$ to compute the risk. 
For every $\lamz \in \{10^{-1}, \cdots, 10^6\}$,
a logistic regression model is trained and evaluated. The $\lamz^*$ value for each distribution is chosen as the value that yields the lowest risk of the algorithm on the test set.

The frequency of each candidate $\lambda$ value being chosen as $\lamz^*$ across the $100$ distributions is shown in Figure~\ref{fig:figure1}. The simulation results, summarized in the frequency plot, show a strong preference for the extreme values of the regularization parameter $\lamz$. This pattern provides empirical evidence supporting the conditional threshold-like behaviour of $\lambda_z$ established in Corollary \ref{cor: bang_bang}.

Importantly, our goal in this simulation is to evaluate the structural properties of the risk surface, not to propose an operational model-selection protocol. Therefore, rather than selecting $\lamz$ via cross-validation on the training set (which would introduce finite-sample model-selection noise), we perform a retrospective scan using the test set risk to approximate the true population risk. While evaluating on the test set is a form of data snooping in an applied setting, here it serves strictly as an exploratory oracle to compute the best-case $\lamz^*$. Because the threshold-like preference relies on the population signal-to-noise ratio $\rho$, this oracle approach confirms that the algorithm empirically exhibits the conditional threshold-like behavior in practice, reflecting the theoretical bound.

\begin{figure}[htbp]
\centering
\includegraphics[width=.42\textwidth]{fig1}
\caption{Frequency of $\lamz^*$ across 100 distributions.}
\label{fig:figure1}
\end{figure}



\paragraph{Real World Data Set:}


To validate the threshold-like behavior in practical settings, we examine the $\lamz^*$ values across eight real-world datasets from the OpenML repository. For each dataset, we identify a primary label $Y$ and a relevant binary side-information variable $Z$, while the remaining covariates are treated as features $X$. The datasets are summarized in Table~\ref{tab:realworld-datasets}.

\begin{table}[htbp]
\centering
\small
\caption{Real-world datasets used in the simulation study.}
\label{tab:realworld-datasets}
\resizebox{\textwidth}{!}{%
\begin{tabular}{l r l l l}
\toprule
\textbf{Dataset} & \textbf{Total Size} & \textbf{$Y$ (Label)} & \textbf{$Z$ (Side Info)} & \textbf{$X$ (Features)} \\
\midrule
Adult Income & $48{,}842$ & Income ($\le$\$50K / $>$\$50K) & Sex (Male / Female) & 13 demographic features \cite{openml:adult} \\
Bank Marketing & $45{,}211$ & Class (1 / 2) & V7 (Yes / No) & 15 client features (V1--V16) \cite{openml:bank} \\
Mushroom & $8{,}124$ & Edibility (Edible / Poisonous) & Bruises (Yes / No) & 21 physical features \cite{openml:mushroom} \\
Spambase & $4{,}601$ & Email Type (Spam / Non-Spam) & Word ``make'' (Pres. / Abs.) & 56 word/capital statistics \cite{openml:spambase} \\
Titanic Survivor & $1{,}309$ & Survival (Survived / Died) & Sex (Male / Female) & 3 features (class, age, fare) \cite{openml:titanic} \\
German Credit & $1{,}000$ & Credit Risk (Good / Bad) & Foreign Worker (Yes / No) & 19 credit/financial features \cite{openml:credit-g} \\
Diabetes & $768$ & Diabetes Status (+ / --) & Age ($\le 30$ / $> 30$) & 7 medical measurement features \cite{openml:diabetes} \\
Breast Cancer & $699$ & Diagnosis (Benign / Malignant) & Clump Thickness ($\le 5$ / $> 5$) & 8 cell measurement features \cite{openml:breast-w} \\
\bottomrule
\end{tabular}%
}
\end{table}

We employ a 70/30 train-test split for each dataset. The training portion is then partitioned using two regimes: an 80/20 split (80\% feature-label, 20\% feature-side information) and a 50/50 split (equal feature-label and feature-side information). Within each regime, for every value of $\lamz \in \{10^{-1}, \cdots, 10^6\}$,
a model is trained and evaluated. The $\lamz^*$ value for each distribution is chosen as the value that yields the lowest risk of the algorithm on the test set.

The results are reported as follows. For all eight datasets, we present a bar plot of the value of $\lambda_z^*$ under both the 80/20 and 50/50 regimes (Figure~\ref{fig:figure2}). The results broadly support the theoretical prediction: in both regimes, most datasets select $\lambda_z^*$ near one of the two extreme values of the candidate grid rather than settling at an intermediate value, with only a few exceptions. 

\begin{figure}[htbp]
\centering
\includegraphics[width=.52\textwidth]{fig2}
\caption{$\lambda_z^*$ across datasets (80/20 and 50/50 regimes).}
\label{fig:figure2}
\end{figure}

Furthermore, for the three largest datasets (Adult, Bank, and Mushroom), we present progressive risk plots (Figure~\ref{fig:figure3}) that illustrate how the test risk evolves as the training pool size increases, under both the 80/20 and 50/50 partitioning regimes. Each panel directly compares three approaches: \emph{Full Side Info} (using both $X$ and $Z$ for prediction), \emph{No Side Info} (using only $X$), and \emph{Partitioned Side Info} (our Algorithm~\ref{alg: side information algorithm}, evaluated at the smallest and largest candidate values of $\lambda_z$). The results in Figure~\ref{fig:figure3} show that, for the partitioned side information algorithm, the minimum risk achieved over the smallest and largest candidate values of $\lamz$ consistently lies between the Full Side Info and No Side Info plots.

\begin{figure}[htbp] 
\centering
\includegraphics[width=.44\textwidth]{fig3}
\caption{Risk plots for largest datasets (80/20 and 50/50 regimes).}
\label{fig:figure3}
\end{figure}

\FloatBarrier

In this section, we present simulation studies comparing the theoretical upper bound on excess risk (Theoretical RHS) with the actual observed test-set excess risk (Actual LHS), each evaluated at the bound-minimizing regularization parameter $\lambda_z^\star$ \coderef{subsec:code_side_exp}. The simulation strictly implements the Partitioned Side Information Learning (PSIL) procedure defined in Algorithm~\ref{alg: side information algorithm}, directly optimizing the regularized objectives for both the hard imputer $\hat{\phi}(x) \in \{-1, +1\}$ and the primary classifier $\hat{w}_\lambda$, and evaluating deployment risk using observed side information $z$ at test time.

To investigate the theoretical guarantees across diverse distributional regimes, we first evaluate three synthetic toy distributions where the covariates $X \in \mathbb{R}^5$ are sampled from a three-component Gaussian mixture model truncated to a compact Euclidean ball to satisfy Assumption (C1). The distributions represent three distinct regimes: highly informative side information, a moderate signal-imputation balance, and noisy side information. For each distribution, we vary the primary sample size $n \in \{500, 1000, 2000, 5000\}$ with auxiliary sample size $m = n$. We calculate the theoretical upper bound evaluated at the optimal candidate $\lambda_z^\star$ using the envelope constants ($C_1 = C_2 = \mathsf{W} = 1$) and compare it with the empirical excess risk achieved on an independent held-out test set of $20{,}000$ samples across $20$ independent repetitions, displaying $99\%$ confidence bands. As shown in Figure~\ref{fig:excess_risk_toy}, the actual excess risk closely tracks the shape and decay rate of the theoretical upper bound across all three regimes.

\begin{figure}[htbp]
\centering
\includegraphics[width=\textwidth]{fig_excess_risk_toy.pdf}
\caption{Theoretical Upper Bound (RHS) and Actual Excess Risk (LHS) with 99\% confidence bands across sample size $n$ ($m=n$) for three truncated Gaussian mixture toy distributions.}
\label{fig:excess_risk_toy}
\end{figure}

We conduct an evaluation on three real-world benchmark datasets from the OpenML repository: Adult Income, Bank Marketing, and Mushroom. To ensure strict separation between model selection and evaluation, each dataset is partitioned into three disjoint splits: a training pool ($70\%$) from which $(S, D)$ are sampled across varying sample sizes $n$ ($m = n$), a separate oracle sample ($15\%$) used exclusively to estimate population parameters and select the candidate regularizer $\lambda_z^\star$, and a held-out test set ($15\%$) used solely to evaluate test-set deployment risk. No test observations influence model selection or regularizer tuning. As displayed in Figure~\ref{fig:excess_risk_real}, the actual excess risk decays steadily with increasing sample size, demonstrating consistent alignment with the trajectory of the theoretical upper bound across diverse tabular feature domains.

\begin{figure}[htbp]
\centering
\includegraphics[width=\textwidth]{fig_excess_risk_real.pdf}
\caption{Theoretical Upper Bound (RHS) and Actual Excess Risk (LHS) with 99\% confidence bands across sample size $n$ ($m=n$) for Adult Income, Bank Marketing, and Mushroom under strict 3-way sample partitioning.}
\label{fig:excess_risk_real}
\end{figure}

\FloatBarrier
